% ── Section: Algorithm ───────────────────────────────────────────────────────
\section{The RBAC-HNSW Algorithm}
\label{sec:algorithm}

\subsection{Bitmask Encoding}
\label{sec:bitmask}

Each vector $v$ is assigned a 64-bit unsigned integer $v_m$ encoding its
access requirements.
A query with permission mask $q_m$ may access $v$ if and only if:
\[
  (v_m \;\mathbin{\&}\; q_m) = q_m
\]
That is, the node must possess \emph{at least} every permission bit required
by the query.
This subsumes standard RBAC role-hierarchy checks: if role $R_1$ inherits
from $R_2$, set all $R_2$-bits inside $R_1$'s mask.
Figure~\ref{fig:rbac_schema} shows the 64-bit layout used in our clinical
benchmark (see Section~\ref{sec:dataset}).

The mask is stored adjacent to the node's vector pointer in memory, so a
single 64-byte cache-line fetch~\cite{drepper2007every} delivers both fields.
The bitwise AND and comparison together constitute one CPU instruction
($\approx\!1$~ns), evaluated \emph{before} the distance computation
($\approx\!30$~ns for 768-dim cosine similarity).

\subsection{Two-Gate Greedy Search}
\label{sec:search}

Algorithm~\ref{alg:rbac_search} presents the modified HNSW layer-0 search.
The key departure from vanilla HNSW is in lines~\ref{alg:gate1}--\ref{alg:route}:
when a neighbor node $n$ fails the RBAC gate, we \emph{do not discard it}.
Instead, we add $n$ to the routing frontier so its own neighbor list can
direct the beam toward accessible regions.
Distance is only computed (line~\ref{alg:gate2}) when the access check passes.

\begin{algorithm}[t]
\small
\caption{RBAC-HNSW Greedy Search (Layer 0)}
\label{alg:rbac_search}
\begin{algorithmic}[1]
\Require Graph $G$, query $\mathbf{q}$, query mask $q_m$, result count $k$, beam $\mathit{ef}$
\Ensure Top-$k$ accessible nearest neighbors
\State $\mathit{results} \leftarrow$ min-heap; $\mathit{routing} \leftarrow$ priority queue; $\mathit{visited} \leftarrow \emptyset$
\State $e \leftarrow \textsc{EntryPoint}(G)$; $\mathit{visited} \leftarrow \{e\}$
\If{$\textsc{RBACCheck}(e, q_m)$} \Comment{Gate 1}
    \State push $(\textsc{Dist}(\mathbf{q}, e),\, e)$ into $\mathit{results}$ and $\mathit{routing}$
\Else
    \State push $(\infty,\, e)$ into $\mathit{routing}$ \Comment{Route through denied}
\EndIf
\While{$\mathit{routing} \neq \emptyset$ \textbf{and} $|\mathit{results}| < \mathit{ef}$}
    \State $(d_c, c) \leftarrow \textsc{Pop}(\mathit{routing})$
    \ForAll{$n \in \textsc{Neighbors}(G, c)$}
        \If{$n \in \mathit{visited}$} \textbf{continue} \EndIf
        \State $\mathit{visited} \leftarrow \mathit{visited} \cup \{n\}$
        \If{$(n.v_m \;\mathbin{\&}\; q_m) \neq q_m$} \label{alg:gate1} \Comment{RBAC gate}
            \State push $(\infty, n)$ into $\mathit{routing}$ \label{alg:route}
            \State \textbf{continue}
        \EndIf
        \State $d \leftarrow \textsc{Dist}(\mathbf{q}, n)$ \label{alg:gate2} \Comment{Distance gate}
        \State push $(d, n)$ into $\mathit{results}$ and $\mathit{routing}$
    \EndFor
\EndWhile
\State \Return $\textsc{TopK}(\mathit{results}, k)$
\end{algorithmic}
\end{algorithm}

\subsection{Why Routing Through Denied Nodes Matters}
\label{sec:routing}

Consider 0.1\% selectivity over one million vectors: roughly 1{,}000 accessible
records scattered throughout a graph of one million nodes.
With standard post-filtering, a beam of width $\mathit{ef} = 200$ explores
200 candidates; the expected number of accessible candidates hit is
$200 \times 0.001 = 0.2$---far below the $k=10$ needed.
Even inflating $\mathit{ef}$ to 10{,}000 would cost 50$\times$ more distance
computations and still struggle with sparse accessible regions.

RBAC-HNSW routes through denied nodes at \emph{zero distance cost} (no
multiply-accumulate).
The beam's exploration radius grows freely through the graph topology until it
reaches an accessible node cluster, then spends its distance budget productively.
This insight is conceptually related to ACORN~\cite{patel2024acorn} (routing
through predicate-failing nodes via augmented neighbor lists) and the relaxed
monotonicity of VBASE~\cite{zhang2023vbase}, though our formulation is
specific to bitmask RBAC and requires no graph augmentation.
Figure~\ref{fig:algorithm} illustrates the traversal.
